% answers_short.tex —— 2015 年·四川卷·理科 · 简短答案速查
% 由 tools/make_short_answers.py 自 answers.tex 的 \gksolhead/\gkanswer 提取，
% 并与 2015/2015_数学真题_四川理_答案要点.md 逐题比对后生成，勿手改（改源 MD 后重跑）。

\begin{gknotice}
\textbf{使用说明}：本节只给最终答案便于快速核对，完整推导见解析版。答案为编写者独立推导，\textbf{非官方标准答案}。
\end{gknotice}

\gkpart{选择题}
\begin{center}\begin{tabular}{ccccccccccc}
\toprule
题号 & 1 & 2 & 3 & 4 & 5 & 6 & 7 & 8 & 9 & 10 \\
\midrule
答案 & A & C & D & A & D & B & C & B & B & D \\
\bottomrule
\end{tabular}\end{center}

\gkpart{填空题}
\begin{itemize}[leftmargin=2.6em]
  \item[\textbf{11.}] \(-40\)
  \item[\textbf{12.}] \(\dfrac{\sqrt{6}}{2}\)
  \item[\textbf{13.}] \(24\) 小时
  \item[\textbf{14.}] \(\dfrac{2}{5}\)
  \item[\textbf{15.}] ①④
\end{itemize}

\gkpart{解答题}
\begin{itemize}[leftmargin=2.6em]
  \item[\textbf{16.}] (1) \(a_n=2^n\)；(2) \(n\) 的最小值为 \(10\)
  \item[\textbf{17.}] (1) \(\dfrac{99}{100}\)；(2) \(X\) 的分布列为 \(P(X=1)=\dfrac{1}{5}\)、\(P(X=2)=\dfrac{3}{5}\)、\(P(X=3)=\dfrac{1}{5}\)，数学期望 \(E(X)=2\)
  \item[\textbf{18.}] (1) \(F\) 标记在 \(B\) 正上方的顶点、\(G\) 标记在 \(C\) 正上方的顶点、\(H\) 标记在 \(D\) 正上方的顶点（使上底面为 \(EFGH\)）；(2) 成立（取 \(P\) 为 \(CD\) 中点，\(MP\parallel BD\)、\(PN\parallel DH\)，故平面 \(MPN\parallel\) 平面 \(BDH\)，从而 \(MN\parallel\) 平面 \(BDH\)）；(3) \(\dfrac{2\sqrt{2}}{3}\)
  \item[\textbf{19.}] (1) 成立（由二倍角公式 \(\sin A=2\sin\dfrac{A}{2}\cos\dfrac{A}{2}\)、\(1-\cos A=2\sin^{2}\dfrac{A}{2}\) 相除即得）；(2) \(\dfrac{4\sqrt{10}}{3}\)
  \item[\textbf{20.}] (1) \(\dfrac{x^2}{4}+\dfrac{y^2}{2}=1\)；(2) 存在，\(Q(0,2)\)
  \item[\textbf{21.}] (1) 分情形——\(a\geqslant\dfrac{1}{4}\) 时 \(g(x)\) 在 \((0,+\infty)\) 上单调递增；\(0<a<\dfrac{1}{4}\) 时 \(g(x)\) 先增、后减、再增（临界点为 \(x^2-x+a=0\) 的两根）；(2) 成立（唯一极小值点 \(x_0(a)\)，\(\Phi(a)=f(x_0(a),a)\) 连续递减，\(\Phi(0^{+})=1>0\)、\(\Phi(1/2)\leqslant f(\mathrm{e},1/2)=\mathrm{e}^2-3\mathrm{e}-1<0\)，由介值定理存在 \(a\in(0,1)\)；结合单调性得 \(f(x)\geqslant 0\) 恒成立且零点唯一）
\end{itemize}
